\chapter{Unit 5: Handling Data with Pandas and Data Visualisation with Matplotlib}

\section{Section A: Multiple Choice Practice Questions (60 Curated Problems)}

\noindent This section contains 60 rigorous multiple-choice questions extracted from university course examination assessments. Review each question, test edge-case syntax, and verify code execution mechanics.
\vspace{3mm}

\noindent\textbf{Q1.} \hfill {\small\color{LPUOrange}\textbf{[introduction to pandas | Easy]}}\\[2pt]
What is Pandas primarily used for in Python?

\begin{enumerate}[label=(\alph*),topsep=2pt,itemsep=1pt,partopsep=0pt,parsep=0pt]
    \item \relax Managing operating system processes
    \item \relax Data manipulation and analysis
    \item \relax Compiling C++ code
    \item \relax Building 3D games
\end{enumerate}
\vspace{2mm}

\noindent\textbf{Q2.} \hfill {\small\color{LPUOrange}\textbf{[introduction to pandas | Easy]}}\\[2pt]
Which is the conventional way to import the Pandas library?

\begin{enumerate}[label=(\alph*),topsep=2pt,itemsep=1pt,partopsep=0pt,parsep=0pt]
    \item \relax \texttt{include pandas as pd}
    \item \relax \texttt{import pd as pandas}
    \item \relax \texttt{import pandas as pd}
    \item \relax \texttt{import pandas.pd}
\end{enumerate}
\vspace{2mm}

\noindent\textbf{Q3.} \hfill {\small\color{LPUOrange}\textbf{[introduction to pandas | Easy]}}\\[2pt]
Pandas is built on top of which numerical computing library?

\begin{enumerate}[label=(\alph*),topsep=2pt,itemsep=1pt,partopsep=0pt,parsep=0pt]
    \item \relax Matplotlib
    \item \relax TensorFlow
    \item \relax SciPy
    \item \relax NumPy
\end{enumerate}
\vspace{2mm}

\noindent\textbf{Q4.} \hfill {\small\color{LPUOrange}\textbf{[series | Easy]}}\\[2pt]
A Pandas Series is best described as a:

\begin{enumerate}[label=(\alph*),topsep=2pt,itemsep=1pt,partopsep=0pt,parsep=0pt]
    \item \relax Two-dimensional labeled table
    \item \relax One-dimensional labeled array
    \item \relax Three-dimensional data cube
    \item \relax Scalar single value
\end{enumerate}
\vspace{2mm}

\noindent\textbf{Q5.} \hfill {\small\color{LPUOrange}\textbf{[series | Easy]}}\\[2pt]
Which function is used to create a Pandas Series?

\begin{enumerate}[label=(\alph*),topsep=2pt,itemsep=1pt,partopsep=0pt,parsep=0pt]
    \item \relax \texttt{pd.DataFrame()}
    \item \relax \texttt{pd.array()}
    \item \relax \texttt{pd.Series()}
    \item \relax \texttt{pd.series()}
\end{enumerate}
\vspace{2mm}

\noindent\textbf{Q6.} \hfill {\small\color{LPUOrange}\textbf{[series | Easy]}}\\[2pt]
What is the default index for a Series created from a list of 4 elements?

\begin{enumerate}[label=(\alph*),topsep=2pt,itemsep=1pt,partopsep=0pt,parsep=0pt]
    \item \relax Integers 1, 2, 3, 4
    \item \relax Integers 0, 1, 2, 3
    \item \relax Letters a, b, c, d
    \item \relax No index is assigned
\end{enumerate}
\vspace{2mm}

\noindent\textbf{Q7.} \hfill {\small\color{LPUOrange}\textbf{[dataframe | Easy]}}\\[2pt]
A Pandas DataFrame is a:

\begin{enumerate}[label=(\alph*),topsep=2pt,itemsep=1pt,partopsep=0pt,parsep=0pt]
    \item \relax One-dimensional labeled array
    \item \relax Single scalar value
    \item \relax Plain Python dictionary
    \item \relax Two-dimensional labeled data structure
\end{enumerate}
\vspace{2mm}

\noindent\textbf{Q8.} \hfill {\small\color{LPUOrange}\textbf{[dataframe | Easy]}}\\[2pt]
Which method displays the first 5 rows of a DataFrame \texttt{df} by default?

\begin{enumerate}[label=(\alph*),topsep=2pt,itemsep=1pt,partopsep=0pt,parsep=0pt]
    \item \relax \texttt{df.first()}
    \item \relax \texttt{df.head()}
    \item \relax \texttt{df.top()}
    \item \relax \texttt{df.start()}
\end{enumerate}
\vspace{2mm}

\noindent\textbf{Q9.} \hfill {\small\color{LPUOrange}\textbf{[dataframe | Easy]}}\\[2pt]
Which attribute gives the number of rows and columns of a DataFrame as a tuple?

\begin{enumerate}[label=(\alph*),topsep=2pt,itemsep=1pt,partopsep=0pt,parsep=0pt]
    \item \relax \texttt{df.dimensions}
    \item \relax \texttt{df.length}
    \item \relax \texttt{df.size}
    \item \relax \texttt{df.shape}
\end{enumerate}
\vspace{2mm}

\noindent\textbf{Q10.} \hfill {\small\color{LPUOrange}\textbf{[working with csv files | Easy]}}\\[2pt]
Which Pandas function reads data from a CSV file into a DataFrame?

\begin{enumerate}[label=(\alph*),topsep=2pt,itemsep=1pt,partopsep=0pt,parsep=0pt]
    \item \relax \texttt{pd.read\_csv()}
    \item \relax \texttt{pd.import\_csv()}
    \item \relax \texttt{pd.load\_csv()}
    \item \relax \texttt{pd.open\_csv()}
\end{enumerate}
\vspace{2mm}

\noindent\textbf{Q11.} \hfill {\small\color{LPUOrange}\textbf{[working with csv files | Easy]}}\\[2pt]
Which method writes a DataFrame \texttt{df} to a CSV file?

\begin{enumerate}[label=(\alph*),topsep=2pt,itemsep=1pt,partopsep=0pt,parsep=0pt]
    \item \relax \texttt{df.export\_csv()}
    \item \relax \texttt{df.write\_csv()}
    \item \relax \texttt{df.save\_csv()}
    \item \relax \texttt{df.to\_csv()}
\end{enumerate}
\vspace{2mm}

\noindent\textbf{Q12.} \hfill {\small\color{LPUOrange}\textbf{[working with csv files | Easy]}}\\[2pt]
What does CSV stand for?

\begin{enumerate}[label=(\alph*),topsep=2pt,itemsep=1pt,partopsep=0pt,parsep=0pt]
    \item \relax Character-Set Values
    \item \relax Common Storage Version
    \item \relax Column-Sorted Values
    \item \relax Comma-Separated Values
\end{enumerate}
\vspace{2mm}

\noindent\textbf{Q13.} \hfill {\small\color{LPUOrange}\textbf{[operations using dataframes | Easy]}}\\[2pt]
Which method provides summary statistics like mean, min, and max for numeric columns?

\begin{enumerate}[label=(\alph*),topsep=2pt,itemsep=1pt,partopsep=0pt,parsep=0pt]
    \item \relax \texttt{df.stats()}
    \item \relax \texttt{df.summary()}
    \item \relax \texttt{df.info()}
    \item \relax \texttt{df.describe()}
\end{enumerate}
\vspace{2mm}

\noindent\textbf{Q14.} \hfill {\small\color{LPUOrange}\textbf{[operations using dataframes | Easy]}}\\[2pt]
How do you select a single column named \texttt{age} from a DataFrame \texttt{df}?

\begin{enumerate}[label=(\alph*),topsep=2pt,itemsep=1pt,partopsep=0pt,parsep=0pt]
    \item \relax \texttt{df['age']}
    \item \relax \texttt{df.column('age')}
    \item \relax \texttt{df.get('age', all=True)}
    \item \relax \texttt{df->age}
\end{enumerate}
\vspace{2mm}

\noindent\textbf{Q15.} \hfill {\small\color{LPUOrange}\textbf{[operations using dataframes | Easy]}}\\[2pt]
Which method is used to remove rows containing missing (NaN) values?

\begin{enumerate}[label=(\alph*),topsep=2pt,itemsep=1pt,partopsep=0pt,parsep=0pt]
    \item \relax \texttt{df.deletena()}
    \item \relax \texttt{df.removeNaN()}
    \item \relax \texttt{df.clearna()}
    \item \relax \texttt{df.dropna()}
\end{enumerate}
\vspace{2mm}

\noindent\textbf{Q16.} \hfill {\small\color{LPUOrange}\textbf{[line plots | Easy]}}\\[2pt]
Which Matplotlib function is used to create a line plot ?

\begin{enumerate}[label=(\alph*),topsep=2pt,itemsep=1pt,partopsep=0pt,parsep=0pt]
    \item \relax \texttt{plt.line()}
    \item \relax \texttt{plt.graph()}
    \item \relax \texttt{plt.plot()}
    \item \relax \texttt{plt.draw()}
\end{enumerate}
\vspace{2mm}

\noindent\textbf{Q17.} \hfill {\small\color{LPUOrange}\textbf{[multiple subplots in one figure | Easy]}}\\[2pt]
Which function creates a figure with multiple subplots at once?

\begin{enumerate}[label=(\alph*),topsep=2pt,itemsep=1pt,partopsep=0pt,parsep=0pt]
    \item \relax \texttt{plt.subplots()}
    \item \relax \texttt{plt.multiplot()}
    \item \relax \texttt{plt.grid()}
    \item \relax \texttt{plt.figures()}
\end{enumerate}
\vspace{2mm}

\noindent\textbf{Q18.} \hfill {\small\color{LPUOrange}\textbf{[histograms | Easy]}}\\[2pt]
A histogram is mainly used to show the:

\begin{enumerate}[label=(\alph*),topsep=2pt,itemsep=1pt,partopsep=0pt,parsep=0pt]
    \item \relax Trend over time only
    \item \relax Distribution of a numeric variable
    \item \relax Proportion of a whole
    \item \relax Relationship between two categories
\end{enumerate}
\vspace{2mm}

\noindent\textbf{Q19.} \hfill {\small\color{LPUOrange}\textbf{[bar charts | Easy]}}\\[2pt]
Which Matplotlib function creates a vertical bar chart ?

\begin{enumerate}[label=(\alph*),topsep=2pt,itemsep=1pt,partopsep=0pt,parsep=0pt]
    \item \relax \texttt{plt.bar()}
    \item \relax \texttt{plt.hist()}
    \item \relax \texttt{plt.barh()}
    \item \relax \texttt{plt.column()}
\end{enumerate}
\vspace{2mm}

\noindent\textbf{Q20.} \hfill {\small\color{LPUOrange}\textbf{[pie charts | Easy]}}\\[2pt]
A pie chart is best suited for showing:

\begin{enumerate}[label=(\alph*),topsep=2pt,itemsep=1pt,partopsep=0pt,parsep=0pt]
    \item \relax Frequency distribution of bins
    \item \relax Proportions of a whole
    \item \relax Correlation between variables
    \item \relax Trends over time
\end{enumerate}
\vspace{2mm}

\noindent\textbf{Q21.} \hfill {\small\color{LPUOrange}\textbf{[introduction to pandas | Medium]}}\\[2pt]
You import pandas using \texttt{import pandas as pd}. Which statement correctly describes why pandas is preferred over plain Python lists for tabular data analysis?

\begin{enumerate}[label=(\alph*),topsep=2pt,itemsep=1pt,partopsep=0pt,parsep=0pt]
    \item \relax It replaces the need for the NumPy library entirely
    \item \relax It compiles Python code into C for faster loops
    \item \relax It provides labeled, aligned data structures with vectorized operations
    \item \relax It stores data only as JSON internally for speed
\end{enumerate}
\vspace{2mm}

\noindent\textbf{Q22.} \hfill {\small\color{LPUOrange}\textbf{[series | Medium]}}\\[2pt]
Given \texttt{s = pd.Series([10, 20, 30], index=['a', 'b', 'c'])}, what does \texttt{s['b']} return?

\begin{enumerate}[label=(\alph*),topsep=2pt,itemsep=1pt,partopsep=0pt,parsep=0pt]
    \item \relax \texttt{10}
    \item \relax \texttt{['b', 20]}
    \item \relax \texttt{30}
    \item \relax \texttt{20}
\end{enumerate}
\vspace{2mm}

\noindent\textbf{Q23.} \hfill {\small\color{LPUOrange}\textbf{[series | Medium]}}\\[2pt]
What is the result of \texttt{pd.Series([1, 2, 3]) + pd.Series([10, 20], index=[0, 2])}?

\begin{enumerate}[label=(\alph*),topsep=2pt,itemsep=1pt,partopsep=0pt,parsep=0pt]
    \item \relax It raises a length mismatch error
    \item \relax All values are added position by position ignoring the index
    \item \relax It concatenates into a Series of length 5
    \item \relax Values at index 1 become \texttt{NaN} due to alignment
\end{enumerate}
\vspace{2mm}

\noindent\textbf{Q24.} \hfill {\small\color{LPUOrange}\textbf{[dataframe | Medium]}}\\[2pt]
For a DataFrame \texttt{df}, which expression selects only the rows where the column \texttt{age} is greater than 30?

\begin{enumerate}[label=(\alph*),topsep=2pt,itemsep=1pt,partopsep=0pt,parsep=0pt]
    \item \relax \texttt{df[df.age.filter(> 30)]}
    \item \relax \texttt{df[df['age'] > 30]}
    \item \relax \texttt{df.loc['age' > 30]}
    \item \relax \texttt{df['age' > 30]}
\end{enumerate}
\vspace{2mm}

\noindent\textbf{Q25.} \hfill {\small\color{LPUOrange}\textbf{[dataframe | Medium]}}\\[2pt]
What does \texttt{df.shape} return for a DataFrame with 100 rows and 5 columns?

\begin{enumerate}[label=(\alph*),topsep=2pt,itemsep=1pt,partopsep=0pt,parsep=0pt]
    \item \relax \texttt{(5, 100)}
    \item \relax \texttt{[100, 5]}
    \item \relax \texttt{(100, 5)}
    \item \relax \texttt{500}
\end{enumerate}
\vspace{2mm}

\noindent\textbf{Q26.} \hfill {\small\color{LPUOrange}\textbf{[dataframe | Medium]}}\\[2pt]
Which method would you use to view the first 5 rows of a DataFrame \texttt{df}?

\begin{enumerate}[label=(\alph*),topsep=2pt,itemsep=1pt,partopsep=0pt,parsep=0pt]
    \item \relax \texttt{df.first(5)}
    \item \relax \texttt{df.top()}
    \item \relax \texttt{df.begin()}
    \item \relax \texttt{df.head()}
\end{enumerate}
\vspace{2mm}

\noindent\textbf{Q27.} \hfill {\small\color{LPUOrange}\textbf{[working with csv files | Medium]}}\\[2pt]
You have a CSV file where columns are separated by semicolons (\texttt{;}). Which call reads it correctly?

\begin{enumerate}[label=(\alph*),topsep=2pt,itemsep=1pt,partopsep=0pt,parsep=0pt]
    \item \relax \texttt{pd.read\_csv('data.csv', delim=';')}
    \item \relax \texttt{pd.read\_csv('data.csv', split=';')}
    \item \relax \texttt{pd.read\_csv('data.csv', sep=';')}
    \item \relax \texttt{pd.load\_csv('data.csv', sep=';')}
\end{enumerate}
\vspace{2mm}

\noindent\textbf{Q28.} \hfill {\small\color{LPUOrange}\textbf{[working with csv files | Medium]}}\\[2pt]
You want to save a DataFrame \texttt{df} to a CSV file without writing the index column. Which call is correct?

\begin{enumerate}[label=(\alph*),topsep=2pt,itemsep=1pt,partopsep=0pt,parsep=0pt]
    \item \relax \texttt{df.write\_csv('out.csv', index=None)}
    \item \relax \texttt{df.to\_csv('out.csv', index=False)}
    \item \relax \texttt{df.save\_csv('out.csv', index=0)}
    \item \relax \texttt{df.to\_csv('out.csv', header=False)}
\end{enumerate}
\vspace{2mm}

\noindent\textbf{Q29.} \hfill {\small\color{LPUOrange}\textbf{[working with csv files | Medium]}}\\[2pt]
When reading a CSV that has no header row, which argument ensures pandas does not treat the first data row as column names?

\begin{enumerate}[label=(\alph*),topsep=2pt,itemsep=1pt,partopsep=0pt,parsep=0pt]
    \item \relax \texttt{skiprows=1}
    \item \relax \texttt{header=None}
    \item \relax \texttt{names=True}
    \item \relax \texttt{header=0}
\end{enumerate}
\vspace{2mm}

\noindent\textbf{Q30.} \hfill {\small\color{LPUOrange}\textbf{[operations using dataframes | Medium]}}\\[2pt]
Given a DataFrame \texttt{df} with a numeric column \texttt{sales}, which expression computes the average of that column?

\begin{enumerate}[label=(\alph*),topsep=2pt,itemsep=1pt,partopsep=0pt,parsep=0pt]
    \item \relax \texttt{df['sales'].mean()}
    \item \relax \texttt{df.mean('sales')}
    \item \relax \texttt{mean(df['sales'])}
    \item \relax \texttt{df['sales'].average()}
\end{enumerate}
\vspace{2mm}

\noindent\textbf{Q31.} \hfill {\small\color{LPUOrange}\textbf{[operations using dataframes | Medium]}}\\[2pt]
What does \texttt{df.groupby('city')['sales'].sum()} produce?

\begin{enumerate}[label=(\alph*),topsep=2pt,itemsep=1pt,partopsep=0pt,parsep=0pt]
    \item \relax A count of rows per city
    \item \relax Total sales for each unique city
    \item \relax The sales sorted alphabetically by city
    \item \relax The overall sum of all sales values
\end{enumerate}
\vspace{2mm}

\noindent\textbf{Q32.} \hfill {\small\color{LPUOrange}\textbf{[operations using dataframes | Medium]}}\\[2pt]
Which method returns the count of missing (\texttt{NaN}) values in each column of a DataFrame \texttt{df}?

\begin{enumerate}[label=(\alph*),topsep=2pt,itemsep=1pt,partopsep=0pt,parsep=0pt]
    \item \relax \texttt{df.isnull().sum()}
    \item \relax \texttt{df.isnull().count()}
    \item \relax \texttt{df.dropna().sum()}
    \item \relax \texttt{df.missing().sum()}
\end{enumerate}
\vspace{2mm}

\noindent\textbf{Q33.} \hfill {\small\color{LPUOrange}\textbf{[operations using dataframes | Medium]}}\\[2pt]
You want to create a new column \texttt{total} equal to \texttt{price} multiplied by \texttt{quantity}. Which statement is correct?

\begin{enumerate}[label=(\alph*),topsep=2pt,itemsep=1pt,partopsep=0pt,parsep=0pt]
    \item \relax \texttt{df.total = multiply(price, quantity)}
    \item \relax \texttt{df['total'] = df['price'].mul()}
    \item \relax \texttt{df.add('total', price * quantity)}
    \item \relax \texttt{df['total'] = df['price'] * df['quantity']}
\end{enumerate}
\vspace{2mm}

\noindent\textbf{Q34.} \hfill {\small\color{LPUOrange}\textbf{[line plots | Medium]}}\\[2pt]
Using matplotlib, which pair of calls plots \texttt{y} against \texttt{x} as a line and then displays it?

\begin{enumerate}[label=(\alph*),topsep=2pt,itemsep=1pt,partopsep=0pt,parsep=0pt]
    \item \relax \texttt{plt.lineplot(x, y); plt.show()}
    \item \relax \texttt{plt.line(x, y); plt.draw()}
    \item \relax \texttt{plt.plot(y, x); plt.render()}
    \item \relax \texttt{plt.plot(x, y); plt.show()}
\end{enumerate}
\vspace{2mm}

\noindent\textbf{Q35.} \hfill {\small\color{LPUOrange}\textbf{[line plots | Medium]}}\\[2pt]
In \texttt{plt.plot(x, y, 'r--')}, what does the format string \texttt{'r--'} specify?

\begin{enumerate}[label=(\alph*),topsep=2pt,itemsep=1pt,partopsep=0pt,parsep=0pt]
    \item \relax A rectangular line style
    \item \relax A red dashed line
    \item \relax Two red solid lines
    \item \relax A round dotted marker
\end{enumerate}
\vspace{2mm}

\noindent\textbf{Q36.} \hfill {\small\color{LPUOrange}\textbf{[multiple subplots in one figure | Medium]}}\\[2pt]
Which call creates a figure with a 2-row by 2-column grid of subplots?

\begin{enumerate}[label=(\alph*),topsep=2pt,itemsep=1pt,partopsep=0pt,parsep=0pt]
    \item \relax \texttt{fig, ax = plt.figure(2, 2)}
    \item \relax \texttt{fig, ax = plt.subplot(2, 2)}
    \item \relax \texttt{fig, ax = plt.grid(2, 2)}
    \item \relax \texttt{fig, ax = plt.subplots(2, 2)}
\end{enumerate}
\vspace{2mm}

\noindent\textbf{Q37.} \hfill {\small\color{LPUOrange}\textbf{[multiple subplots in one figure | Medium]}}\\[2pt]
After \texttt{fig, axes = plt.subplots(2, 2)}, how do you draw on the subplot in the bottom-right corner?

\begin{enumerate}[label=(\alph*),topsep=2pt,itemsep=1pt,partopsep=0pt,parsep=0pt]
    \item \relax \texttt{axes['bottom-right'].plot(...)}
    \item \relax \texttt{axes(1, 1).plot(...)}
    \item \relax \texttt{axes[2, 2].plot(...)}
    \item \relax \texttt{axes[1, 1].plot(...)}
\end{enumerate}
\vspace{2mm}

\noindent\textbf{Q38.} \hfill {\small\color{LPUOrange}\textbf{[histograms | Medium]}}\\[2pt]
Which parameter of \texttt{plt.hist(data, bins=20)} controls the number of intervals the data range is divided into?

\begin{enumerate}[label=(\alph*),topsep=2pt,itemsep=1pt,partopsep=0pt,parsep=0pt]
    \item \relax \texttt{width}
    \item \relax \texttt{bins}
    \item \relax \texttt{range}
    \item \relax \texttt{intervals}
\end{enumerate}
\vspace{2mm}

\noindent\textbf{Q39.} \hfill {\small\color{LPUOrange}\textbf{[histograms | Medium]}}\\[2pt]
A histogram is most appropriate for visualizing which of the following?

\begin{enumerate}[label=(\alph*),topsep=2pt,itemsep=1pt,partopsep=0pt,parsep=0pt]
    \item \relax The proportion of parts in a whole
    \item \relax The relationship between two categorical variables
    \item \relax The frequency distribution of a single continuous variable
    \item \relax The exact value of each individual record
\end{enumerate}
\vspace{2mm}

\noindent\textbf{Q40.} \hfill {\small\color{LPUOrange}\textbf{[bar charts | Medium]}}\\[2pt]
You have categories in a list \texttt{cats} and their counts in \texttt{vals}. Which call creates a vertical bar chart?

\begin{enumerate}[label=(\alph*),topsep=2pt,itemsep=1pt,partopsep=0pt,parsep=0pt]
    \item \relax \texttt{plt.plot(cats, vals, 'bar')}
    \item \relax \texttt{plt.bar(cats, vals)}
    \item \relax \texttt{plt.hist(cats, vals)}
    \item \relax \texttt{plt.barh(cats, vals)}
\end{enumerate}
\vspace{2mm}

\noindent\textbf{Q41.} \hfill {\small\color{LPUOrange}\textbf{[series | Hard]}}\\[2pt]
Given \texttt{s = pd.Series([10, 20, 30], index=['a', 'b', 'c'])}, what does \texttt{s[0:2]} return compared to \texttt{s['a':'b']}?

\begin{enumerate}[label=(\alph*),topsep=2pt,itemsep=1pt,partopsep=0pt,parsep=0pt]
    \item \relax \texttt{s[0:2]} raises a \texttt{KeyError} since integer indexing is disabled when labels exist
    \item \relax Both return exactly the same two elements because slicing is always position-based
    \item \relax \texttt{s['a':'b']} excludes 'b' because all Python slicing is exclusive of the endpoint
    \item \relax \texttt{s[0:2]} returns elements at positions 0 and 1; \texttt{s['a':'b']} returns elements 'a' and 'b' (label slicing is inclusive of endpoint)
\end{enumerate}
\vspace{2mm}

\noindent\textbf{Q42.} \hfill {\small\color{LPUOrange}\textbf{[dataframe | Hard]}}\\[2pt]
For a DataFrame \texttt{df}, which statement correctly distinguishes \texttt{df.loc[]} from \texttt{df.iloc[]} when the index is \texttt{[2, 0, 1]}?

\begin{enumerate}[label=(\alph*),topsep=2pt,itemsep=1pt,partopsep=0pt,parsep=0pt]
    \item \relax Both select the same row because index labels are ignored
    \item \relax \texttt{df.loc[0]} raises an error because the index is not sorted
    \item \relax \texttt{df.iloc[0]} selects the row labeled 0 and \texttt{df.loc[0]} selects the first physical row
    \item \relax \texttt{df.loc[0]} selects the row with index label 0; \texttt{df.iloc[0]} selects the first physical row (label 2)
\end{enumerate}
\vspace{2mm}

\noindent\textbf{Q43.} \hfill {\small\color{LPUOrange}\textbf{[operations using dataframes | Hard]}}\\[2pt]
Given \texttt{df1} with index \texttt{['a','b']} and \texttt{df2} with index \texttt{['b','c']}, what is the result of \texttt{df1 + df2} for overlapping and non-overlapping labels?

\begin{enumerate}[label=(\alph*),topsep=2pt,itemsep=1pt,partopsep=0pt,parsep=0pt]
    \item \relax All rows are summed positionally regardless of labels
    \item \relax Rows 'a' and 'c' become \texttt{NaN} (no alignment match); row 'b' holds the element-wise sum
    \item \relax The operation raises a \texttt{ValueError} due to mismatched indices
    \item \relax Only row 'b' survives; rows 'a' and 'c' are dropped from the result
\end{enumerate}
\vspace{2mm}

\noindent\textbf{Q44.} \hfill {\small\color{LPUOrange}\textbf{[working with csv files | Hard]}}\\[2pt]
A CSV has a column of integers but some cells are empty. After \texttt{pd.read\_csv('data.csv')}, what dtype does that column typically have and why?

\begin{enumerate}[label=(\alph*),topsep=2pt,itemsep=1pt,partopsep=0pt,parsep=0pt]
    \item \relax \texttt{int64}, because pandas fills missing integers with 0 automatically
    \item \relax \texttt{Int64}, because pandas always uses nullable integer types by default
    \item \relax \texttt{float64}, because \texttt{NaN} (a float) forces the entire integer column to be upcast to float
    \item \relax \texttt{object}, because any missing value converts all numbers to strings
\end{enumerate}
\vspace{2mm}

\noindent\textbf{Q45.} \hfill {\small\color{LPUOrange}\textbf{[operations using dataframes | Hard]}}\\[2pt]
For a numeric DataFrame \texttt{df}, how do \texttt{df.apply(np.sum, axis=0)} and \texttt{df.apply(np.sum, axis=1)} differ?

\begin{enumerate}[label=(\alph*),topsep=2pt,itemsep=1pt,partopsep=0pt,parsep=0pt]
    \item \relax Both return a single scalar equal to the total sum of all elements
    \item \relax \texttt{axis=0} returns per-row sums; \texttt{axis=1} returns per-column sums
    \item \relax \texttt{axis=0} returns one sum per column (collapsing rows); \texttt{axis=1} returns one sum per row (collapsing columns)
    \item \relax \texttt{axis=1} raises an error because \texttt{np.sum} only works column-wise
\end{enumerate}
\vspace{2mm}

\noindent\textbf{Q46.} \hfill {\small\color{LPUOrange}\textbf{[dataframe | Hard]}}\\[2pt]
What is the key behavioral difference between \texttt{df.dropna(how='all')} and \texttt{df.dropna(how='any')} (default)?

\begin{enumerate}[label=(\alph*),topsep=2pt,itemsep=1pt,partopsep=0pt,parsep=0pt]
    \item \relax \texttt{how='all'} fills \texttt{NaN} with the mean instead of dropping
    \item \relax Both drop identical rows; the parameter only changes performance
    \item \relax \texttt{how='all'} drops rows with any \texttt{NaN}; \texttt{how='any'} drops only fully-\texttt{NaN} rows
    \item \relax \texttt{how='all'} drops a row only if every value is \texttt{NaN}; \texttt{how='any'} drops a row if at least one value is \texttt{NaN}
\end{enumerate}
\vspace{2mm}

\noindent\textbf{Q47.} \hfill {\small\color{LPUOrange}\textbf{[line plots | Hard]}}\\[2pt]
When plotting \texttt{plt.plot(x, y1); plt.plot(x, y2)} on the same axes without specifying colors, what happens?

\begin{enumerate}[label=(\alph*),topsep=2pt,itemsep=1pt,partopsep=0pt,parsep=0pt]
    \item \relax The second \texttt{plot} call overwrites and erases the first line
    \item \relax Both lines are drawn in the same color, making them indistinguishable
    \item \relax A \texttt{ValueError} is raised because colors must be specified explicitly
    \item \relax Matplotlib automatically assigns different colors from the default color cycle to each line
\end{enumerate}
\vspace{2mm}

\noindent\textbf{Q48.} \hfill {\small\color{LPUOrange}\textbf{[multiple subplots in one figure | Hard]}}\\[2pt]
Using \texttt{fig, ax = plt.subplots(2, 2)}, how should you correctly access the subplot in the bottom-right position?

\begin{enumerate}[label=(\alph*),topsep=2pt,itemsep=1pt,partopsep=0pt,parsep=0pt]
    \item \relax \texttt{ax[2][2]}, using 1-based row and column numbers
    \item \relax \texttt{ax[3]}, because subplots are stored in a flat 1D list
    \item \relax \texttt{ax[1, 1]}, because \texttt{ax} is a 2D NumPy array of Axes indexed by [row, column]
    \item \relax \texttt{ax.bottom\_right}, using the named-position attribute
\end{enumerate}
\vspace{2mm}

\noindent\textbf{Q49.} \hfill {\small\color{LPUOrange}\textbf{[histograms | Hard]}}\\[2pt]
For \texttt{plt.hist(data, bins=10)} on data ranging from 0 to 100, what determines each bin's width and count?

\begin{enumerate}[label=(\alph*),topsep=2pt,itemsep=1pt,partopsep=0pt,parsep=0pt]
    \item \relax Width is fixed at 1 unit regardless of the data range
    \item \relax Width is the range divided by bins (10 units each); count is the number of data points falling in each interval
    \item \relax The count is the cumulative total of all points up to that bin's upper edge
    \item \relax Each bin always contains exactly the same number of data points (equal-frequency binning)
\end{enumerate}
\vspace{2mm}

\noindent\textbf{Q50.} \hfill {\small\color{LPUOrange}\textbf{[bar charts | Hard]}}\\[2pt]
What is the fundamental difference between \texttt{plt.bar()} and \texttt{plt.hist()}?

\begin{enumerate}[label=(\alph*),topsep=2pt,itemsep=1pt,partopsep=0pt,parsep=0pt]
    \item \relax \texttt{bar()} can only produce horizontal bars while \texttt{hist()} produces vertical ones
    \item \relax \texttt{bar()} bins continuous data automatically while \texttt{hist()} requires pre-counted categories
    \item \relax \texttt{bar()} plots one bar per categorical/discrete value you supply; \texttt{hist()} bins continuous data and plots frequencies of those bins
    \item \relax Both are identical; \texttt{hist} is just an alias for \texttt{bar}
\end{enumerate}
\vspace{2mm}

\noindent\textbf{Q51.} \hfill {\small\color{LPUOrange}\textbf{[pie charts | Hard]}}\\[2pt]
In \texttt{plt.pie(sizes, autopct='\%1.1f\%\%')}, what does \texttt{autopct} control and how are wedge sizes determined?

\begin{enumerate}[label=(\alph*),topsep=2pt,itemsep=1pt,partopsep=0pt,parsep=0pt]
    \item \relax \texttt{autopct} controls wedge color; sizes are determined by insertion order only
    \item \relax \texttt{autopct} sets the number of wedges; wedge sizes are always equal
    \item \relax \texttt{autopct} rotates the chart; wedge sizes come from raw values without normalization
    \item \relax \texttt{autopct} formats each wedge's percentage label; wedge angles are proportional to each value divided by the total
\end{enumerate}
\vspace{2mm}

\noindent\textbf{Q52.} \hfill {\small\color{LPUOrange}\textbf{[introduction to pandas | Hard]}}\\[2pt]
Why can a pandas \texttt{Series} hold heterogeneous data more flexibly than a plain NumPy array of a fixed numeric dtype?

\begin{enumerate}[label=(\alph*),topsep=2pt,itemsep=1pt,partopsep=0pt,parsep=0pt]
    \item \relax A \texttt{Series} stores each element in a separate NumPy array of its own dtype
    \item \relax A \texttt{Series} can fall back to the \texttt{object} dtype, storing Python objects of mixed types, though at the cost of vectorized performance
    \item \relax A \texttt{Series} internally converts all values to strings to allow mixing types
    \item \relax A \texttt{Series} cannot hold mixed types; it always raises an error like NumPy
\end{enumerate}
\vspace{2mm}

\noindent\textbf{Q53.} \hfill {\small\color{LPUOrange}\textbf{[operations using dataframes | Hard]}}\\[2pt]
Given \texttt{df.groupby('cat')['val'].transform('mean')}, how does its output differ from \texttt{df.groupby('cat')['val'].mean()}?

\begin{enumerate}[label=(\alph*),topsep=2pt,itemsep=1pt,partopsep=0pt,parsep=0pt]
    \item \relax \texttt{transform} drops the grouping while \texttt{mean} keeps every original row
    \item \relax \texttt{transform} returns a Series aligned to the original rows (each row gets its group mean); \texttt{mean} returns one aggregated value per group
    \item \relax \texttt{transform} returns a scalar; \texttt{mean} returns a full DataFrame
    \item \relax Both return one value per group with identical shape
\end{enumerate}
\vspace{2mm}

\noindent\textbf{Q54.} \hfill {\small\color{LPUOrange}\textbf{[working with csv files | Hard]}}\\[2pt]
What is the effect of \texttt{pd.read\_csv('data.csv', index\_col=0)} versus omitting \texttt{index\_col}?

\begin{enumerate}[label=(\alph*),topsep=2pt,itemsep=1pt,partopsep=0pt,parsep=0pt]
    \item \relax Both produce identical results because the first column always becomes the index
    \item \relax \texttt{index\_col=0} skips the first column entirely rather than using it as an index
    \item \relax With \texttt{index\_col=0}, the first column becomes the DataFrame index; without it, a default \texttt{RangeIndex} (0,1,2,...) is created and all columns are kept as data
    \item \relax Omitting \texttt{index\_col} causes the last column to be used as the index by default
\end{enumerate}
\vspace{2mm}

\noindent\textbf{Q55.} \hfill {\small\color{LPUOrange}\textbf{[dataframe | Hard]}}\\[2pt]
Which comparison of \texttt{df['col']} and \texttt{df[['col']]} is correct?

\begin{enumerate}[label=(\alph*),topsep=2pt,itemsep=1pt,partopsep=0pt,parsep=0pt]
    \item \relax \texttt{df[['col']]} raises an error because double brackets are invalid syntax
    \item \relax \texttt{df['col']} returns a 1D Series; \texttt{df[['col']]} returns a single-column DataFrame (2D)
    \item \relax \texttt{df['col']} returns a DataFrame; \texttt{df[['col']]} returns a Series
    \item \relax Both return a Series with identical shape and type
\end{enumerate}
\vspace{2mm}

\noindent\textbf{Q56.} \hfill {\small\color{LPUOrange}\textbf{[multiple subplots in one figure | Hard]}}\\[2pt]
When \texttt{fig, axes = plt.subplots(1, 3, sharey=True)} is used, what does \texttt{sharey=True} accomplish?

\begin{enumerate}[label=(\alph*),topsep=2pt,itemsep=1pt,partopsep=0pt,parsep=0pt]
    \item \relax It forces all subplots to share the same x-axis instead of the y-axis
    \item \relax All three subplots share the same y-axis scale and limits, and inner y-tick labels are hidden for clarity
    \item \relax The y-data of the three plots is summed into a single shared plot
    \item \relax Each subplot gets an independent y-axis automatically scaled to its own data
\end{enumerate}
\vspace{2mm}

\noindent\textbf{Q57.} \hfill {\small\color{LPUOrange}\textbf{[series | Hard]}}\\[2pt]
For \texttt{s = pd.Series([1, 2, 3])}, why does \texttt{s + pd.Series([1, 2, 3], index=[1, 2, 3])} produce \texttt{NaN} values?

\begin{enumerate}[label=(\alph*),topsep=2pt,itemsep=1pt,partopsep=0pt,parsep=0pt]
    \item \relax Labels are aligned before addition; only indices 1 and 2 overlap, so indices 0 and 3 yield \texttt{NaN}
    \item \relax The operation fails entirely and returns an empty Series
    \item \relax Addition is positional, so all three pairs add cleanly with no \texttt{NaN}
    \item \relax \texttt{NaN} appears because the two Series have different lengths
\end{enumerate}
\vspace{2mm}

\noindent\textbf{Q58.} \hfill {\small\color{LPUOrange}\textbf{[histograms | Hard]}}\\[2pt]
In \texttt{plt.hist(data, bins=20, density=True)}, what does \texttt{density=True} change about the y-axis?

\begin{enumerate}[label=(\alph*),topsep=2pt,itemsep=1pt,partopsep=0pt,parsep=0pt]
    \item \relax The y-axis shows cumulative counts up to each bin
    \item \relax The y-axis shows raw counts multiplied by the number of bins
    \item \relax The y-axis shows probability density so the total area under the bars equals 1, rather than raw counts
    \item \relax The y-axis is unchanged; \texttt{density} only affects bar color
\end{enumerate}
\vspace{2mm}

\noindent\textbf{Q59.} \hfill {\small\color{LPUOrange}\textbf{[line plots | Hard]}}\\[2pt]
Given time-series \texttt{x} values that are unsorted, what visual artifact does \texttt{plt.plot(x, y)} produce and why?

\begin{enumerate}[label=(\alph*),topsep=2pt,itemsep=1pt,partopsep=0pt,parsep=0pt]
    \item \relax It raises an error demanding sorted x-values
    \item \relax Matplotlib automatically sorts x before drawing, so the line is always smooth
    \item \relax Only the points are shown with no connecting line at all
    \item \relax The line zig-zags back and forth because \texttt{plot} connects points in the order given, not in sorted x-order
\end{enumerate}
\vspace{2mm}

\noindent\textbf{Q60.} \hfill {\small\color{LPUOrange}\textbf{[operations using dataframes | Hard]}}\\[2pt]
What does \texttt{df.merge(other, on='key', how='outer')} produce compared to \texttt{how='inner'}?

\begin{enumerate}[label=(\alph*),topsep=2pt,itemsep=1pt,partopsep=0pt,parsep=0pt]
    \item \relax Both keep the same rows; \texttt{how} only reorders columns
    \item \relax \texttt{outer} keeps only matching keys; \texttt{inner} keeps everything
    \item \relax \texttt{outer} concatenates rows vertically without matching on keys
    \item \relax \texttt{outer} keeps all keys from both frames (filling unmatched sides with \texttt{NaN}); \texttt{inner} keeps only keys present in both
\end{enumerate}
\vspace{2mm}


\newpage
\section{Section B: Comprehensive Subjective Problem with Detailed Analytical Rubric}

\begin{subjectiveq}{5}{1}
Construct an end-to-end data processing and visualization pipeline with Pandas and Matplotlib:
\begin{enumerate}[label=(\alph*)]
    \item Load a sales dataset, impute missing numeric values with column medians, and remove duplicate rows based on composite keys.
    \item Compute monthly aggregated revenue, moving averages, and filter customer segments whose cumulative revenue exceeds the 90th percentile.
    \item Generate a multi-panel Matplotlib publication figure displaying: (1) Time series line plot with confidence bands, (2) Distribution histogram with KDE curve, and (3) Categorical horizontal bar chart.
\end{enumerate}
\end{subjectiveq}

\begin{solbox}
\textbf{Pandas \& Matplotlib Pipeline:}
\begin{pycode}
import pandas as pd\\
import matplotlib.pyplot as plt\\
~\\
def clean\_and\_aggregate(df):\\
\quad df = df.drop\_duplicates(subset=['transaction\_id', 'customer\_id']).copy()\\
\quad num\_cols = df.select\_dtypes(include=['float64', 'int64']).columns\\
\quad df[num\_cols] = df[num\_cols].fillna(df[num\_cols].median())\\
\quad df['date'] = pd.to\_datetime(df['date'])\\
\quad monthly = df.resample('M', on='date')['revenue'].sum().reset\_index()\\
\quad monthly['rolling\_3m'] = monthly['revenue'].rolling(window=3).mean()\\
\quad return df, monthly\\
~\\
def plot\_dashboard(monthly, df):\\
\quad fig, axes = plt.subplots(1, 2, figsize=(14, 5))\\
\quad axes[0].plot(monthly['date'], monthly['revenue'], label='Monthly Revenue', marker='o')\\
\quad axes[0].plot(monthly['date'], monthly['rolling\_3m'], label='3M MA', linestyle='--')\\
\quad axes[0].set\_title("Revenue Dynamics")\\
\quad axes[0].legend()\\
\quad axes[1].hist(df['revenue'], bins=30, color='skyblue', edgecolor='black')\\
\quad axes[1].set\_title("Revenue Distribution")\\
\quad plt.tight\_layout()\\
\quad plt.savefig('sales\_dashboard.png', dpi=300)
\end{pycode}
\textbf{Marking Insights:} Emphasize correct use of \texttt{.loc}, explicit copies to avoid \texttt{SettingWithCopyWarning}, and memory optimization through categorical conversions.
\end{solbox}

\begin{examinertip}
\begin{itemize}[leftmargin=1.5em, itemsep=2pt]
    \item \textbf{Algorithmic Formulation (3 Marks):} Correct problem decomposition, function signatures, and parameter documentation.
    \item \textbf{Syntactic Rigor \& Implementation (4 Marks):} Valid Python syntax, proper error handling, defensive input validation, and avoidance of anti-patterns.
    \item \textbf{Complexity \& Edge Analysis (3 Marks):} Accurate time and space asymptotic analysis $\mathcal{O}(n)$ and testing edge conditions (null inputs, zero division, empty collections).
\end{itemize}
\end{examinertip}
